% ===================================================================
%  圖表第 3 號 — 童年同青年。
%
%  Traduction, faite en 2026, en cantonais ecrit d'aujourd'hui, en
%  caracteres traditionnels, sur l'ido de texto/io. Regles de la
%  colonne : voir 10-toubiu-01.tex. Deux scenes, deux numerotations :
%  les cles portent c1 et c2.
%
%  LE CORTEGE DE BAPTEME OBLIGE LE CANTONAIS A CE QUE LE VIETNAMIEN
%  AVAIT DEJA DEMANDE : nommer chaque parent par sa place. Mais il le
%  fait autrement, et plus finement encore, parce qu'il compte l'AGE
%  RELATIF a l'interieur d'une meme generation :
%
%    伯父 / 叔父   le frere aine / cadet du pere
%    姑媽 / 姑姐   la soeur ainee / cadette du pere
%    舅父          le frere de la mere
%    姨媽 / 姨姐   la soeur ainee / cadette de la mere
%    表哥 / 表弟   le cousin aine / cadet
%    家姐 / 細妹   la grande soeur / la petite soeur
%
%  L'ido dit « onklulo » et « onklino » sans plus ; la gravure, elle,
%  montre un couple d'age mur derriere le parrain, et le texte les
%  place avant les cousins. On a donc lu la planche et ecrit 伯父 et
%  伯娘, les aines du cote paternel. La petite Berthe est tenue par sa
%  表姐, sa cousine ainee, et son frere marche derriere : 佢阿哥.
%
%  LA FETE DE LA DEUXIEME SCENE EST CELLE QUI FAIT ENTRER LES MOTS DU
%  PORT DE HONG KONG : 沙展 le sergent (de l'anglais), 士多啤梨 la
%  fraise, 巴士 l'autobus n'y sont pas — mais 波 la balle du tableau
%  2 revient, et « 咪高峰 » n'a pas encore lieu d'etre. Ce qui entre
%  ici, ce sont les noms de la foire : 馬騮戲 le cirque des singes,
%  嘉年華 la fete foraine, 過山車 les montagnes russes.
%
%  ET LE MAT DE COCAGNE PORTE ICI SON VRAI NOM. Le vietnamien avait
%  du forger un compose ; le cantonais dit « 爬竿 », grimper au mat,
%  parce que le jeu existe dans les fetes de village du Guangdong.
% ===================================================================

%%K t03-apar-1 apar
\VUsaut{3.0mm}
\VUtitre{15.0pt}{60}{圖表第 3 號}
\VUsaut{1.6mm}
\VUtitre{11.4pt}{0}{童年同青年。}
\VUsaut{3.4mm}

%%K t03-apar-2 apar
（第 3 同第 4 張圖表咁樣配埋一齊，係為咗喺四個\VUgras{家庭場面}入面，講出關於
\VUgras{天氣}、\VUgras{年紀}、\VUgras{家庭}、\VUgras{衣著}同\VUgras{身分}嘅基本
概念。）

\VUsaut{4.0mm}
%%K t03-tit-1 sub
\VUcentre{12.0pt}{0}{\textit{第一場。}}
\VUsaut{3.0mm}

%%K t03-tit-2 sub
\VUpk{10.2pt}{40}{鄉下嘅洗禮。}
\VUsaut{2.4mm}

%%K t03-tit-3 sub
\VUpk{10.2pt}{40}{春天。}
\VUsaut{3.0mm}

%%K t03-c1-01-1 p
1. --- 我哋而家係\VUgras{春天}，喺\VUgras{三月}、\VUgras{四月}或者\VUgras{五月}，
喺\VUgras{一個星期}入面嘅\VUgras{星期一}、\VUgras{星期二}或者\VUgras{星期三}。
而家係朝早\VUgras{十點}。

%%K t03-c1-02-1 p
2. --- \VUgras{天氣}好好，\VUgras{空氣}清新。日頭照住。幾朵\VUgras{雲仔}
\textsuperscript{(2)} 喺藍色嘅\VUgras{天} \textsuperscript{(1)} 上面升起。

%%K t03-c1-03-1 p
3. --- 啲\VUgras{樹}啱啱先至有\VUgras{葉}。幼身嘅\VUgras{白楊} \textsuperscript{(3)}
同高大嘅\VUgras{梧桐} \textsuperscript{(17)} 到而家仲淨係有芽。

%%K t03-c1-04-1 p
4. --- 啲\VUgras{燕} \textsuperscript{(4)} 返嚟喇，佢哋開開心心咁吱吱喳喳，繞住個
\VUgras{尖頂} \textsuperscript{(7)} 飛——嗰個係\VUgras{鐘樓} \textsuperscript{(5)} 嘅
尖頂，而鐘樓就企喺村裏面間\VUgras{教堂} \textsuperscript{(6)} 上面。

%%K t03-tit-4 sub
\VUsaut{3.4mm}
\VUpk{10.2pt}{40}{教堂。}
\VUsaut{3.0mm}

%%K t03-c1-05-1 p
5. --- 呢間教堂好簡樸，有度大\VUgras{門} \textsuperscript{(13)} 同一個大
\VUgras{鐘} \textsuperscript{(8)}。\VUgras{短針} \textsuperscript{(9)} 指住 X 字，
佢指\VUgras{鐘數}。\VUgras{長針} \textsuperscript{(10)} 指住 XII（半個字）；佢指
\VUgras{分鐘}。

%%K t03-c1-06-1 p
6. --- 鐘樓入面，啲\VUgras{鐘} \textsuperscript{(11)} 開開心心咁響。屋頂尖上企住
一個細細嘅石\VUgras{十字架} \textsuperscript{(12)}。

%%K t03-c1-07-1 p
7. --- 教堂唔淨止有度大\VUgras{門} \textsuperscript{(13)}，佢仲有度側門仔。
\VUgras{神父} \textsuperscript{(15)} 就係由呢度側門，著住\VUgras{長袍}行出
\VUgras{聖器房} \textsuperscript{(14)}，去\VUgras{神父樓} \textsuperscript{(19)}。

%%K t03-c1-08-1 p
8. --- 喺佢附近，就係\VUgras{賣牛奶嘅婦人} \textsuperscript{(24)} 同佢隻\VUgras{驢}
\textsuperscript{(23)}。佢由城入面返嚟，喺嗰度佢送咗好牛奶畀啲細路。

%%K t03-tit-5 sub
\VUsaut{4.0mm}
\VUpk{10.2pt}{40}{果園。}
\VUsaut{3.4mm}

%%K t03-c1-09-1 p
9. --- 阿黎伯先生（隻驢）對今朝呢趟嘢好滿意。佢好陶醉咁吸住滿係\VUgras{花}
\textsuperscript{(20)} 香嘅空氣，嗰啲花開滿咗隔籬\VUgras{果園}裏面嘅\VUgras{果樹}
\textsuperscript{(18)}。嗰啲\VUgras{桃樹} \textsuperscript{(18)} 同\VUgras{杏樹} \textsuperscript{(19)}
一片粉紅雪白，睇嚟會有好收成。

%%K t03-c1-10-1 p
10. --- 同時，細細隻嘅\VUgras{蜜蜂} \textsuperscript{(22)} 由\VUgras{蜂竇}
\textsuperscript{(21)} 嗡嗡聲飛去嗰度，採佢哋甜嘅\VUgras{蜜糖}。

%%K t03-c1-11-1 p
11. --- 不過發生咩事呢？原來係村裏面啲細路走埋嚟，嚇到啲\VUgras{鵝}
\textsuperscript{(25)} 走晒。原來係公證人個女\VUgras{領洗}之後，隊人由教堂行出嚟。

%%K t03-c1-12-1 p
12. --- 細細個嘅雅各喺自己\VUgras{搖籃} \textsuperscript{(26)} 入面，畀開心嘅鐘聲
嘈醒；佢家姐瑪德蓮都想睇領洗嘅隊伍，就求佢阿媽出嚟喺門檻度同佢梳頭著衫。

%%K t03-c1-13-1 p
13. --- 阿媽應承咗。佢戴上自己嘅\VUgras{頭巾} \textsuperscript{(28)}，著上
\VUgras{胸衣} \textsuperscript{(29)} 同\VUgras{圍裙} \textsuperscript{(30)}，然後出嚟
坐喺門前一張細凳度。瑪德蓮身上仲淨係得件\VUgras{襯裙} \textsuperscript{(31)} 同件
\VUgras{束腰} \textsuperscript{(32)}。

%%K t03-c1-14-1 p
14. --- 佢幾咁開心呀！佢連自己至愛嘅花都唔記得咗：佢啲\VUgras{石竹}
\textsuperscript{(34)}，有\VUgras{蝴蝶} \textsuperscript{(35)} 停喺上面；佢啲
\VUgras{鬱金香} \textsuperscript{(36)}；佢啲\VUgras{鈴蘭} \textsuperscript{(37)}，
裝喺\VUgras{花盆} \textsuperscript{(38)} 入面，擺喺\VUgras{牆} \textsuperscript{(33)}
上面；仲有砌成咁靚一行籬笆嘅\VUgras{玫瑰} \textsuperscript{(39)}。佢望住隊伍裏面啲太太
嘅華麗衫裙。

%%K t03-c1-15-1 p
15. --- 村裏面啲細路唔多理啲衫。佢哋衝埋去，你推我撞，想攞到\VUgras{糖}
\textsuperscript{(55)} 或者\VUgras{散銀} \textsuperscript{(52)}。其中一個喊咗
\textsuperscript{(40)}。

%%K t03-c1-16-1 p
16. --- 另一個踩親隻\VUgras{狗} \textsuperscript{(42)} 隻腳，隻狗嗌住走，又嚇到隻
\VUgras{雞乸} \textsuperscript{(43)} 同佢啲\VUgras{雞仔} \textsuperscript{(44)} 走晒。

%%K t03-tit-6 sub
\VUsaut{5.0mm}
\VUpk{10.2pt}{40}{領洗嘅隊伍。}
\VUsaut{4.0mm}

%%K t03-c1-17-1 p
17. --- 行喺隊頭嘅係\VUgras{奶媽} \textsuperscript{(45)}。佢戴住一頂靚
\VUgras{布帽} \textsuperscript{(46)}，有長長嘅絲帶。佢抱住個\VUgras{初生嬰兒}
\textsuperscript{(47)}，個嬰兒成身裹住\VUgras{花邊} \textsuperscript{(48)}。

%%K t03-c1-18-1 p
18. --- 奶媽後面係\VUgras{代父} \textsuperscript{(49)} 同\VUgras{代母}
\textsuperscript{(53)}。佢哋派糖同散銀。代父戴住\VUgras{手套} \textsuperscript{(51)}
同一頂\VUgras{高禮帽} \textsuperscript{(50)}。代母手上攞住一束靚\VUgras{花}
\textsuperscript{(54)}。

%%K t03-c1-19-1 p
19. --- 佢哋後面，我哋見到個細路嘅\VUgras{伯父} \textsuperscript{(56)} 同
\VUgras{伯娘} \textsuperscript{(58)}。伯娘戴住一頂好雅致嘅\VUgras{帽}
\textsuperscript{(59)}，伯父就著住黑色\VUgras{禮服大衣} \textsuperscript{(57)} 同
白背心。跟住係\VUgras{表哥} \textsuperscript{(60)} 同\VUgras{表姐}
\textsuperscript{(61)}。表姐拖住初生嬰兒個\VUgras{家姐} \textsuperscript{(62)}，
細細個嘅貝爾達。

%%K t03-c1-20-1 p
20. --- 貝爾達另一個\VUgras{阿哥} \textsuperscript{(63)} 行喺後面。佢拖住一個
\VUgras{親戚} \textsuperscript{(64)}。屋企嘅\VUgras{朋友} \textsuperscript{(65)}
跟住嚟。

%%K t03-tit-7 sub
\VUsaut{4.6mm}
\VUpk{10.2pt}{40}{睇熱鬧嘅人。}
\VUsaut{3.4mm}

%%K t03-c1-21-1 p
21. --- 隊伍附近有個\VUgras{乞兒} \textsuperscript{(66)}，扶住佢支\VUgras{柺杖}
\textsuperscript{(67)}。佢喺度討錢（佢乞食）。

%%K t03-c1-22-1 p
22. --- 另一邊有個好\VUgras{有禮貌} \textsuperscript{(68)} 嘅細路仔。佢向自己識得
嘅人打招呼，講「\VUgras{早晨}」。

%%K t03-c1-23-1 p
23. --- 村裏面啲人行晒出屋，睇領洗嘅隊伍行過。我哋見到一個戴住\VUgras{棉帽}嘅
\VUgras{鄉下佬} \textsuperscript{(69)}，隔籬係個\VUgras{鄉下婆} \textsuperscript{(70)}。
跟住係個\VUgras{老婆婆} \textsuperscript{(71)}，扶住佢支\VUgras{枴杖}
\textsuperscript{(72)}。最後係\VUgras{磨坊佬} \textsuperscript{(73)}，戴住頂大
\VUgras{氈帽} \textsuperscript{(74)}，仲有個著住\VUgras{木屐} \textsuperscript{(75)}
嘅後生鄉下女，
佢喺度教個細路仔行路。

%%K t03-c1-24-1 p
24. --- 呢個場面好得意。

\VUsaut{4.0mm}
%%K t03-tit-8 sub
\VUcentre{12.0pt}{0}{\textit{第二場。}}
\VUsaut{3.0mm}

%%K t03-tit-9 sub
\VUpk{10.2pt}{40}{公眾嘅節日。}
\VUsaut{2.4mm}

%%K t03-tit-10 sub
\VUpk{10.2pt}{40}{水上遊戲同賽艇。}
\VUsaut{3.0mm}

%%K t03-c2-01-1 p
1. --- \VUgras{夏天}到喇。\VUgras{六月}（七月或者\VUgras{八月}）火辣辣嘅日頭引到
啲後生仔去游水同划艇。

%%K t03-c2-02-1 p
2. --- 喺大河右岸，啲划艇嘅人喺度除衫，準備參加水上遊戲同賽艇。

%%K t03-c2-03-1 p
3. --- 其中一個除緊自己對\VUgras{鞋} \textsuperscript{(1)}；佢解緊鞋帶（解鈕）。
佢兩個同伴已經準備好。佢哋而家身上淨係得件\VUgras{冷衫} \textsuperscript{(2)} 同條
\VUgras{泳褲} \textsuperscript{(3)}。佢哋一邊等朋友一邊傾偈。佢哋講緊對岸嘅市集同
節日。

%%K t03-c2-04-1 p
4. --- 佢哋隔籬地下放住同伴嘅\VUgras{衫}：一\VUgras{對}\VUgras{短靴}
\textsuperscript{(4)}、幾對\VUgras{鞋} \textsuperscript{(5)}、\VUgras{襪}
\textsuperscript{(6)}、幾頂\VUgras{氈帽} \textsuperscript{(7)} 同\VUgras{草帽}
\textsuperscript{(8)}，仲有一條\VUgras{領呔} \textsuperscript{(9)}。

%%K t03-c2-05-1 p
5. --- 其中一個後生仔好小心咁摺好自己件\VUgras{外套} \textsuperscript{(10)}。佢將
佢放喺一條毛巾上面，佢隔籬嗰個一陣就會用呢條毛巾將兩人嘅\VUgras{衫}包埋。

%%K t03-c2-06-1 p
6. --- 第六個男仔遲咗。佢頭上仲戴住頂\VUgras{鴨舌帽} \textsuperscript{(11)}。佢仲
未除\VUgras{恤衫} \textsuperscript{(12)}，未除\VUgras{硬領} \textsuperscript{(13)}，
亦未除\VUgras{袖口} \textsuperscript{(14)}。佢手上攞住件\VUgras{背心}
\textsuperscript{(17)}，身上仲著住條\VUgras{長褲} \textsuperscript{(16)} 同吊住條
\VUgras{褲吊帶} \textsuperscript{(15)}。佢實趕唔切下一場比賽。

%%K t03-c2-07-1 p
7. --- 喺啲後生仔附近，有條兩邊生滿\VUgras{薊} \textsuperscript{(18)} 嘅小路通去
河邊，喺嗰度雅各伯喺\VUgras{蘆葦} \textsuperscript{(19)} 叢入面準備緊夜晚要放嘅
\VUgras{煙花} \textsuperscript{(20)}。

%%K t03-tit-11 sub
\VUsaut{4.4mm}
\VUpk{10.2pt}{40}{比賽。}
\VUsaut{3.4mm}

%%K t03-c2-08-1 p
8. --- 河上面幾位太太打住\VUgras{遮}坐喺\VUgras{艇} \textsuperscript{(21)} 度遊河。
佢哋睇緊嗰場好精彩嘅比賽。每一隻\VUgras{賽艇} \textsuperscript{(22)} 上面，四個
\VUgras{划手} \textsuperscript{(24)} 用盡力划，而\VUgras{舵手} \textsuperscript{(26)}
就揸住個\VUgras{舵} \textsuperscript{(25)}，控制住隻薄薄嘅艇。啲\VUgras{槳}
\textsuperscript{(23)} 合住拍子拍水，輕身嘅艇飛咁快咁走。

%%K t03-c2-09-1 p
9. --- 幾個後生仔喺橋躉附近爭\VUgras{水上滑竿} \textsuperscript{(28)} 嘅彩。其中
一個小心咁行喺又軟又滑嘅竿上面，想攞到竿尾嗰面細細嘅\VUgras{旗}
\textsuperscript{(29)}。佢個倒霉嘅\VUgras{對手} \textsuperscript{(30)} 就跌咗落水。
佢游返上岸。

%%K t03-c2-10-1 p
10. --- 大啲嘅後生仔喺度\VUgras{水上比武} \textsuperscript{(32)}，就喺\VUgras{碼頭}
\textsuperscript{(33)} 附近。好多\VUgras{觀眾} \textsuperscript{(31)} 靠喺\VUgras{橋}
\textsuperscript{(27)} 嘅\VUgras{欄杆}度、擠喺\VUgras{河岸}度睇呢啲玩意。不過有幾個
睇熱鬧嘅人回轉頭去睇\VUgras{爬竿} \textsuperscript{(45)} 嘅遊戲，又或者去望
\VUgras{市長嘅隊伍}行嚟\VUgras{派}獎。

%%K t03-c2-11-1 p
11. --- 頭先係幾個\VUgras{頑童} \textsuperscript{(35)} 行喺啲\VUgras{樂手}
\textsuperscript{(36)} 前面；跟住嚟嘅係\VUgras{市長} \textsuperscript{(37)}、佢啲
助理同小城嘅\VUgras{市議會}。行最後嘅係\VUgras{消防員} \textsuperscript{(38)}。

%%K t03-tit-12 sub
\VUsaut{5.0mm}
\VUpk{10.2pt}{40}{市集。}
\VUsaut{3.6mm}

%%K t03-c2-12-1 p
12. --- \VUgras{市集廣場} \textsuperscript{(39)} 上面人山人海。人哋嚟睇唔同嘅
\VUgras{棚}：\VUgras{木馬} \textsuperscript{(40)}、\VUgras{馬戲班}
\textsuperscript{(41)}——嗰度已經開場——仲有\VUgras{公眾舞場} \textsuperscript{(42)}，
城裏面啲後生男女喺嗰度盡情\VUgras{跳舞}。

%%K t03-c2-13-1 p
13. --- 入面盡頭見到\VUgras{鬥牛場} \textsuperscript{(44)}，勇敢嘅西班牙
\VUgras{鬥牛士}喺嗰度同暴怒嘅\VUgras{公牛} \textsuperscript{(45)} 相鬥；跟住係
\VUgras{過山車} \textsuperscript{(49)} 同\VUgras{打靶棚} \textsuperscript{(46)}，
人哋喺嗰度學用\VUgras{槍}打\VUgras{靶}。

%%K t03-c2-14-1 p
14. --- 就喺附近係\VUgras{江湖戲棚} \textsuperscript{(47)}，喺嗰度做\VUgras{喜劇}
同\VUgras{正劇}。緊挨住嘅係\VUgras{巨人} \textsuperscript{(48)} 同\VUgras{侏儒}嘅
棚。前者嘅\VUgras{身高}係兩米一十五；後者就只不過同個細路一樣咁高。

%%K t03-c2-15-1 p
15. --- 最後就係大\VUgras{動物棚} \textsuperscript{(50)}，入面有啲\VUgras{猛獸}：
爪好利嘅\VUgras{老虎} \textsuperscript{(51)}、鬃毛好密嘅\VUgras{獅子}
\textsuperscript{(52)}、兩隻\VUgras{長頸鹿} \textsuperscript{(53)}，仲有隻用
\VUgras{鼻}用得好靈嘅\VUgras{大笨象} \textsuperscript{(54)}。

%%K t03-c2-16-1 p
16. --- 訓練呢啲野獸嘅\VUgras{馴獸師} \textsuperscript{(55)} 企喺棚門口。
\VUsaut{6.0mm}
\VUfilet{22mm}
